\SixthArticle{11}
Wollten wir alle Beziehungen zwischen je drei Funktionen $F(\alpha,\beta,\gamma)$, $F(\alpha+\lambda,\beta+\mu,\gamma+\nu)$, $F(\alpha+\lambda',\beta+\mu',\gamma+\nu')$, in denen $\lambda,\mu,\nu,\lambda',\mu',\nu'$ entweder $=0$ oder $=+1$ oder $=-1$ sind, erschöpfend aufstellen, so würde die Anzahl der Formeln auf 325 steigen. Eine solche Sammlung, wenigstens der einfacheren dieser Formeln, wäre nicht nutzlos. Hier möge es jedoch genügen, nur einige angegeben zu haben, die jeder ohne Mühe entweder aus den Formeln des Artikels 7 oder, wenn er es vorzieht, auf ähnliche Weise beweisen kann, wie die ersten beiden von jenen in Artikel 8 hergeleitet wurden.
\[
\begin{aligned}&[16]\quad F(\alpha,\beta,\gamma)-F(\alpha,\beta,\gamma-1)\\&{}=-\frac{\alpha\beta x}{\gamma(\gamma-1)}F(\alpha+1,\beta+1,\gamma+1)
\end{aligned}
\]
\[
\begin{aligned}&[17]\quad F(\alpha,\beta+1,\gamma)-F(\alpha,\beta,\gamma)\\&{}=\frac{\alpha x}{\gamma}F(\alpha+1,\beta+1,\gamma+1)
\end{aligned}
\]
\[
\begin{aligned}&[18]\quad F(\alpha+1,\beta,\gamma)-F(\alpha,\beta,\gamma)\\&{}=\frac{\beta x}{\gamma}F(\alpha+1,\beta+1,\gamma+1)
\end{aligned}
\]
\[
\begin{aligned}&[19]\quad F(\alpha,\beta+1,\gamma+1)-F(\alpha,\beta,\gamma)\\&{}=\frac{\alpha(\gamma-\beta)x}{\gamma(\gamma+1)}F(\alpha+1,\beta+1,\gamma+2)
\end{aligned}
\]
\[
\begin{aligned}&[20]\quad F(\alpha+1,\beta,\gamma+1)-F(\alpha,\beta,\gamma)\\&{}=\frac{\beta(\gamma-\alpha)x}{\gamma(\gamma+1)}F(\alpha+1,\beta+1,\gamma+2)
\end{aligned}
\]
\[
\begin{aligned}&[21]\quad F(\alpha-1,\beta+1,\gamma)-F(\alpha,\beta,\gamma)\\&{}=\frac{(\alpha-\beta-1)x}{\gamma}F(\alpha,\beta+1,\gamma+1)
\end{aligned}
\]
\[
\begin{aligned}&[22]\quad F(\alpha+1,\beta-1,\gamma)-F(\alpha,\beta,\gamma)\\&{}=\frac{(\beta-\alpha-1)x}{\gamma}F(\alpha+1,\beta,\gamma+1)
\end{aligned}
\]
\[
\begin{aligned}&[23]\quad F(\alpha-1,\beta+1,\gamma)-F(\alpha+1,\beta-1,\gamma)\\&{}=\frac{(\alpha-\beta)x}{\gamma}F(\alpha+1,\beta+1,\gamma+1)
\end{aligned}
\]
\begin{center}\rule{.24\linewidth}{.3pt}\end{center}
\SourceLeaf{143}
\par\Needspace{9\baselineskip}\pdfbookmark[1]{Kettenbrüche}{sixth-section-2}\begin{center}\textbf{\textit{ZWEITER ABSCHNITT.}}\\[6pt]\textbf{\textit{Kettenbrüche.}}\end{center}
\SixthArticle{12}
Bezeichnet man
\[
\frac{F(\alpha,\beta+1,\gamma+1,x)}{F(\alpha,\beta,\gamma,x)}\quad\hbox{mit}\quad G(\alpha,\beta,\gamma,x)
\]
so ergibt sich
\[
\frac{F(\alpha+1,\beta,\gamma+1,x)}{F(\alpha,\beta,\gamma,x)}=\frac{F(\beta,\alpha+1,\gamma+1,x)}{F(\beta,\alpha,\gamma,x)}=G(\beta,\alpha,\gamma,x)
\]
und daher durch Division der Gleichung 19 durch $F(\alpha,\beta+1,\gamma+1,x)$
\[
1-\frac{1}{G(\alpha,\beta,\gamma,x)}=\frac{\alpha(\gamma-\beta)}{\gamma(\gamma+1)}x\,G(\beta+1,\alpha,\gamma+1,x)
\]
oder
\[
[24]\qquad G(\alpha,\beta,\gamma,x)=\cfrac{1}{1-\frac{\alpha(\gamma-\beta)}{\gamma(\gamma+1)}x\,G(\beta+1,\alpha,\gamma+1,x)}
\]
und da ebenso
\[
G(\beta+1,\alpha,\gamma+1,x)=\cfrac{1}{1-\frac{(\beta+1)(\gamma+1-\alpha)}{(\gamma+1)(\gamma+2)}x\,G(\alpha+1,\beta+1,\gamma+2,x)}
\]
usw. gilt, ergibt sich für $G(\alpha,\beta,\gamma,x)$ der Kettenbruch
\[
[25]\qquad\frac{F(\alpha,\beta+1,\gamma+1,x)}{F(\alpha,\beta,\gamma,x)}=\cfrac{1}{1-\cfrac{ax}{1-\cfrac{bx}{1-\cfrac{cx}{1-\cfrac{dx}{1-\hbox{usw.}}}}}}
\]
wobei
\[
\begin{array}{ll}
a=\frac{\alpha(\gamma-\beta)}{\gamma(\gamma+1)}&b=\frac{(\beta+1)(\gamma+1-\alpha)}{(\gamma+1)(\gamma+2)}\\[6pt]
c=\frac{(\alpha+1)(\gamma+1-\beta)}{(\gamma+2)(\gamma+3)}&d=\frac{(\beta+2)(\gamma+2-\alpha)}{(\gamma+3)(\gamma+4)}\\[6pt]
e=\frac{(\alpha+2)(\gamma+2-\beta)}{(\gamma+4)(\gamma+5)}&f=\frac{(\beta+3)(\gamma+3-\alpha)}{(\gamma+5)(\gamma+6)}
\end{array}
\]
usw.; das Fortsetzungsgesetz ist offensichtlich.

Ferner folgt aus den Gleichungen 17, 18, 21, 22
\SourceLeaf{144}
\[
\begin{array}{rl}
[26]&\frac{F(\alpha,\beta+1,\gamma,x)}{F(\alpha,\beta,\gamma,x)}=\cfrac{1}{1-\frac{\alpha x}{\gamma}G(\beta+1,\alpha,\gamma,x)}\\[12pt]
[27]&\frac{F(\alpha+1,\beta,\gamma,x)}{F(\alpha,\beta,\gamma,x)}=\cfrac{1}{1-\frac{\beta x}{\gamma}G(\alpha+1,\beta,\gamma,x)}\\[12pt]
[28]&\frac{F(\alpha-1,\beta+1,\gamma,x)}{F(\alpha,\beta,\gamma,x)}=\cfrac{1}{1-\frac{(\alpha-\beta-1)x}{\gamma}G(\beta+1,\alpha-1,\gamma,x)}\\[12pt]
[29]&\frac{F(\alpha+1,\beta-1,\gamma,x)}{F(\alpha,\beta,\gamma,x)}=\cfrac{1}{1-\frac{(\beta-\alpha-1)x}{\gamma}G(\alpha+1,\beta-1,\gamma,x)}
\end{array}
\]
Setzt man hierin für die Funktion $G$ ihre Werte als Kettenbrüche ein, so ergeben sich ebenso viele neue Kettenbrüche.

Im Übrigen ist von selbst klar, dass der Kettenbruch in Formel 25 abbricht, wenn eine der Zahlen $\alpha,\beta,\gamma-\alpha,\gamma-\beta$ eine negative ganze Zahl ist, andernfalls aber ins Unendliche fortläuft.
\SixthArticle{13}
Die im vorigen Artikel hergeleiteten Kettenbrüche sind von größter Bedeutung, und man kann behaupten, dass die Analysten bisher kaum Kettenbrüche mit einem offensichtlichen Fortsetzungsgesetz gefunden haben, die nicht als Sonderfälle in unseren enthalten sind. Besonders bemerkenswert ist der Fall, in dem in Formel 25 $\beta=0$ gesetzt wird. Daraus folgt $F(\alpha,\beta,\gamma,x)=1$ und somit, wenn man $\gamma-1$ für $\gamma$ schreibt,
\[
\begin{aligned}
{}[30]\quad F(\alpha,1,\gamma)&=1+\frac{\alpha}{\gamma}x+\frac{\alpha(\alpha+1)}{\gamma(\gamma+1)}xx\\&{}+\frac{\alpha(\alpha+1)(\alpha+2)}{\gamma(\gamma+1)(\gamma+2)}x^3+\hbox{usw.}\\[9pt]
&=\cfrac{1}{1-\cfrac{ax}{1-\cfrac{bx}{1-\cfrac{cx}{1-\cfrac{dx}{1-\hbox{usw.}}}}}}
\end{aligned}
\]
wobei
\[
\begin{array}{ll}
a=\frac{\alpha}{\gamma}&b=\frac{\gamma-\alpha}{\gamma(\gamma+1)}\\[6pt]
c=\frac{(\alpha+1)\gamma}{(\gamma+1)(\gamma+2)}&d=\frac{2(\gamma+1-\alpha)}{(\gamma+2)(\gamma+3)}\\[6pt]
e=\frac{(\alpha+2)(\gamma+1)}{(\gamma+3)(\gamma+4)}&f=\frac{3(\gamma+2-\alpha)}{(\gamma+4)(\gamma+5)}
\end{array}
\]
\begin{center}usw.\end{center}
